% Kurzbeschreibung Matching.py
% Maximal 30 Zeilen, Stichpunkte / kurze Sätze

\textbf{Matching.py -- Job zu Kandidaten}

Lädt eine ausgewählte Stelle aus der Tabelle \texttt{jobs} anhand der Job-ID.
Normalisiert Positions-, Fachbereichs- und Ortsangaben der Stelle.
Lädt alle aktiven Kandidaten aus der Tabelle \texttt{candidates} (Status \glqq not interested\grqq{} wird ausgeschlossen).
Bereinigt und normalisiert Kandidatenfelder wie \texttt{position\_now}, \texttt{department}, \texttt{wohnort} und \texttt{wunscharbeitsort}.
Interpretiert \texttt{regionale\_verfuegbarkeit}: \glqq deutschlandweit\grqq{} wird als sehr großer Radius (1000 km) gesetzt, sonst werden Kilometerangaben aus dem Text extrahiert.
Ermittelt für Stellen- und Kandidatenorte geografische Koordinaten aus der Datei \texttt{Staedte\_Deutschland.csv}.
Berechnet Distanzen (Fahrtweg) zwischen Job-Ort und Wohn- bzw. Wunscharbeitsort (inkl. Bundesland- und DE-Sonderfällen).
Nutzt Karrierestufen-Logik (\texttt{KARRIERE\_PFADE}), um zu prüfen, ob die aktuelle Position des Kandidaten zur Zielposition passt.
Prüft, ob der Fachbereich des Kandidaten zum Fachbereich der Stelle passt.
Bewertet die Übereinstimmung von Gehaltswunsch und angebotener Gehaltsspanne mit einem diskreten Gehalts-Score (1--5).
Bewertet den Fahrtweg relativ zur regionalen Verfügbarkeit mit einem Fahrtweg-Score (1--5).
Optional: wertet Skills aus den Job-Anforderungen gegen Kandidaten-Skills aus.
Kandidaten mit passenden Skills erhalten einen Skill-Score; Kandidaten ohne hinterlegte Skills werden nicht ausgeschlossen, erhalten aber keinen Skill-Score.
Kandidaten mit vorhandenen, aber vollständig unpassenden Skills werden vom Matching ausgeschlossen.
Kombiniert Gehalts-, Fahrtweg- und Skill-Score zu einem Gesamt-Score (gewichtetes arithmetisches Mittel).
Berechnet zusätzlich eine Kennzahl zur Datenvollständigkeit (Anteil der vorhandenen Score-Komponenten).
Liefert eine nach Score und Datenvollständigkeit sortierte Kandidatenliste zurück.
Kann die Matching-Ergebnisse inkl. vollständiger Kandidaten- und Job-Daten in eine Excel-Datei exportieren.

